% 03-methodology.tex
\chapter{Methodik}
\label{chap:methodology}

Beschreibe, wie du vorgegangen bist: Experimente, theoretische Vorüberlegungen, Evaluation von Modellen (z.\,B. \ac{AI}).

\section{Forschungsdesign}
\begin{figure}[h]
  \centering
  \includegraphics[width=0.6\linewidth]{example-image-a} % Ersatzgrafik — ersetze durch echte Datei
  \caption{Übersicht über den Forschungsprozess}
  \label{fig:research-process}
\end{figure}

\section{Datenbeschaffung}
\begin{itemize}
  \item Datengrundlage: \dots
  \item Rechtliche Rahmenbedingungen: \dots
\end{itemize}

\section{Mathematische Modelle}
Für die Auswertung der Daten wird häufig auf mathematische Modelle zurückgegriffen. Die Standardabweichung $\sigma$ lässt sich beispielsweise wie in Gleichung~\ref{eq:std_dev} berechnen:
\begin{equation}
  \sigma = \sqrt{\frac{1}{N} \sum_{i=1}^{N} (x_i - \mu)^2}
  \label{eq:std_dev}
\end{equation}
Hierbei beschreibt $N$ die Anzahl der Messwerte, $x_i$ den jeweiligen Messwert und $\mu$ den Mittelwert der Population.

\section{Datenvorbereitung (Code-Beispiel)}
Um die eingelesenen Daten zu verarbeiten, wurde ein Skript in Python erstellt. Listing~\ref{lst:csv_loader} zeigt, wie externe Code-Dateien direkt in \LaTeX{} eingebunden werden können, ohne sie abtippen zu müssen.

\lstinputlisting[
  language=Python,
  caption={Einlesen und Visualisieren von CSV-Daten mittels \texttt{pandas}.},
  label={lst:csv_loader}
]{code/csv_loader.py}

\clearpage